\section{Unidad 3}

\subsection{¿Los portadores mayoritarios que se difunden de la unión polarizada en directa contribuirán con la corriente de base $I_B$ o si pasarán directamente al material tipo $p$?}

Como el material tipo $n$ es muy delgado y su conductividad es baja, un número muy pequeño de estos portadores tomarán esta ruta de alta resistencia hacia la base. El mayor número de estos portadores mayoritarios se difundirá a través de la unión polarizada en inversa hacia el material tipo $n$ conectado al colector.

\subsection{¿Qué relación encuentra entre el símbolo y las polarizaciones de las uniones?}

\subsection{¿Cómo podemos describir lo que sucede en las regiones de corte y saturación?}

En la región de corte las uniones base-emisor y colector-base de un transistor se polarizan en inversa.  
En la región de saturación las uniones base-emisor y colector-base se polarizan en directa.

\subsection{¿Qué nivel de $I_C$ no debería ser superado?}

\[
    (20V) I_C = 300\,\text{mW}
\]
\[I_C = \dfrac{300\,\text{mW}}{20\,\text{V}} = 15\,\text{mA}
\]
